\fancyhead[L]{\small Beweis eines algebraischen Lehrsatzes}
\SourceLeaf{074}
\begin{center}
\textsc{Beweis}\par\medskip
{\large\bfseries EINES ALGEBRAISCHEN LEHRSATZES.}\par\medskip
\rule{0.28\linewidth}{0.4pt}
\end{center}
Gegenstand dieses Aufsatzes ist der \textsc{Cartes}sche, gewöhnlich nach \textsc{Harriot} benannte Satz über den Zusammenhang zwischen der Anzahl der positiven und negativen Wurzeln einer algebraischen Gleichung und der Anzahl der Wechsel und Folgen in den Vorzeichen ihrer Koeffizienten. Den von verschiedenen Autoren versuchten Beweisen dieses Satzes fehlen die Klarheit, Kürze und umfassende Allgemeinheit, die man bei einem so elementaren Gegenstand mit Recht verlangen kann. Eine neue Behandlung scheint daher nicht überflüssig zu sein.

Es sei $X$ ein Polynom in $x$ vom Grad $m$, geordnet nach absteigenden Potenzen von $x$. Wir nehmen ohne Einschränkung der Allgemeinheit an, dass das höchste Glied $x^m$ ist und das konstante, von $x$ unabhängige Glied nicht fehlt. Nur die tatsächlich vorhandenen Glieder sollen aufgeschrieben werden; etwa fehlende Glieder werden also nicht mit dem Koeffizienten $0$ eingefügt.

Sind nicht alle Koeffizienten positiv, so weisen sie einen oder mehrere Vorzeichenwechsel auf. Es sei $-Nx^n$ das erste negative Glied, $+Px^p$ das erste darauf folgende positive und $-Qx^q$ das erste darauf folgende negative Glied usw. Somit sind $m,n,p,q$ usw. abnehmende ganze Zahlen und $N,P,Q$ usw. positiv, und $X$ hat folgende Darstellung:
\[
X=x^m+\,+\, .\, .\,-Nx^n-\,-\, .\, .\,+Px^p+\,+\, .\, .\,-Qx^q-\text{usw.}
\]
\SourceLeaf{075}
Man multipliziere $X$ mit dem linearen Faktor $x-\alpha$, wobei $\alpha$ als positiv vorausgesetzt wird. Man sieht leicht, dass im Produkt $x^{n+1}$ einen negativen, $x^{p+1}$ einen positiven und $x^{q+1}$ einen negativen Koeffizienten hat usw., sodass das Produkt die Form
\[
X(x-\alpha)=x^{m+1}\ldots-N'x^{n+1}\ldots
+P'x^{p+1}\, .\, .\,-Q'x^{q+1}\ldots
\]
annimmt, wobei $N',P',Q'$ usw. positiv sind. Die Vorzeichen zwischen den aufgeführten Gliedern bleiben zwar unbestimmt. Doch ist klar, dass vom höchsten Glied bis zur Potenz $x^{n+1}$ mindestens ein Vorzeichenwechsel, bis $x^{p+1}$ mindestens zwei, bis $x^{q+1}$ mindestens drei usw. auftreten. Liegt der letzte Vorzeichenwechsel in $X$ beim Glied $\pm Ux^u$ und bezeichnet man den Koeffizienten von $x^{u+1}$ in $X(x-\alpha)$ mit $\pm U'$, so ist $U'$ positiv, und bis zum Glied $\pm U'x^{u+1}$ sind mindestens ebenso viele Vorzeichenwechsel aufgetreten, wie es in $X$ gibt. Das letzte Glied in $X(x-\alpha)$ hat aber das Vorzeichen $\mp$; bis dahin muss also mindestens ein weiterer Vorzeichenwechsel hinzugekommen sein. Wir schließen somit, dass $X(x-\alpha)$ mindestens einen Vorzeichenwechsel mehr hat als $X$.

Es sei nun $X$ das Produkt aller linearen Faktoren, die den negativen und imaginären Wurzeln einer Gleichung $y=0$ entsprechen. Sind also $\alpha,\beta,\gamma$ usw. die positiven Wurzeln derselben Gleichung, so ist
\[
y=X(x-\alpha)(x-\beta)(x-\gamma)\ldots
\]
Nach dem vorstehenden Satz hat somit $X(x-\alpha)$ mindestens einen Vorzeichenwechsel mehr als $X$, $X(x-\alpha)(x-\beta)$ mindestens zwei mehr, $X(x-\alpha)(x-\beta)(x-\gamma)$ mindestens drei mehr usw. Folglich gibt es in $y$, selbst wenn in $X$ gar kein Vorzeichenwechsel vorkommt, mindestens ebenso viele Vorzeichenwechsel wie positive Wurzeln. Es ist unmittelbar klar, dass man, wenn die Gleichung weder negative noch imaginäre Wurzeln hat, $X=1$ setzen muss und dieser Schluss gültig bleibt.

Es gehe $y$ in $y'$ über, wenn man die Vorzeichen der Koeffizienten der Potenzen $x^{m-1}$, $x^{m-3}$, $x^{m-5}$ usw. umkehrt. Sämtliche Wurzeln der Gleichung $y'=0$ sind dann den Wurzeln der Gleichung $y=0$ entgegengesetzt. Daher gibt es in $y'$ mindestens ebenso viele Vorzeichenwechsel, wie die Gleichung $y=0$ negative Wurzeln hat.
\par
\SourceLeaf{076}
Wir haben daher folgenden Satz:

\emph{Die Gleichung $y=0$ kann nicht mehr positive Wurzeln haben, als es Vorzeichenwechsel in $y$ gibt, und nicht mehr negative Wurzeln, als es Vorzeichenwechsel in $y'$ gibt.}

Diese Formulierung des Satzes scheint die zweckmäßigste zu sein, da sie die größte Einfachheit mit der umfassendsten Allgemeinheit verbindet und alle Formen des Satzes, die nur unter besonderen Bedingungen gelten, von selbst aus ihr folgen.

Will man die Schranke für die Anzahl der negativen Wurzeln unmittelbar an den Vorzeichen der Koeffizienten von $y$ ablesen, so muss man die \emph{unmittelbaren} Vorzeichenwechsel und Vorzeichenfolgen, bei Gliedern, deren Exponenten von $x$ sich um eins unterscheiden, von den durch fehlende Glieder \emph{unterbrochenen} unterscheiden. Offenbar entspricht jeder unmittelbare und jeder durch eine gerade Anzahl fehlender Glieder unterbrochene Vorzeichenwechsel in $y'$ einer gleichartigen Vorzeichenfolge in $y$, während ein durch eine ungerade Anzahl fehlender Glieder unterbrochener Vorzeichenwechsel in $y'$ auch in $y$ ein gleichartiger Vorzeichenwechsel bleibt. Der zweite Teil des Satzes lässt sich daher auch so ausdrücken:

Die Anzahl der negativen Wurzeln der Gleichung $y=0$ kann nicht größer sein als die Anzahl der unmittelbaren und der durch eine gerade Anzahl fehlender Glieder unterbrochenen Vorzeichenfolgen, vermehrt um die Anzahl der durch eine ungerade Anzahl fehlender Glieder unterbrochenen Vorzeichenwechsel in $y$.

Fehlt in $y$ kein Glied, so ist die Anzahl der negativen Wurzeln nicht größer als die Anzahl der Vorzeichenfolgen.

Bezeichnet $A$ die Anzahl der unmittelbaren Vorzeichenwechsel und $B$ die Anzahl der unmittelbaren Vorzeichenfolgen in $y$, so ist, wenn kein Glied fehlt, $A+B=m$, also gleich der Anzahl aller Wurzeln. Da diese Vorzeichen nur besagen, dass die Anzahl der positiven Wurzeln nicht größer als $A$ und die der negativen nicht größer als $B$ sein kann, bleibt unentschieden, ob und wie viele imaginäre Wurzeln vorhanden sind. Weiß man aber aus anderer Quelle, dass die Gleichung keine imaginären. Wurzeln hat, so muss $A$ notwendig der Anzahl der positiven und $B$ der Anzahl der negativen Wurzeln gleich sein.

Anders verhält es sich, wenn in $y$ Glieder fehlen. Um klar
\SourceLeaf{077}
zu übersehen, was sich daraus über die imaginären Wurzeln schließen lässt, bezeichnen wir mit $a$ die Anzahl der durch eine gerade und mit $c$ die Anzahl der durch eine ungerade Anzahl fehlender Glieder unterbrochenen Vorzeichenwechsel; mit $b$ und $d$ jeweils die Anzahl der durch eine gerade und eine ungerade Anzahl fehlender Glieder unterbrochenen Vorzeichenfolgen in $y$. Man sieht leicht, dass $m-A-B-a-b-c-d$ gleich der Anzahl sämtlicher fehlender Glieder ist, die wir mit $e$ bezeichnen wollen. Nach unserem Satz ist nun die Anzahl der positiven Wurzeln höchstens $A+a+c$ und die Anzahl der negativen höchstens $B+b+c$, also die Anzahl aller reellen Wurzeln höchstens
\[
A+B+a+b+2c=m+c-d-e
\]
Die Anzahl der imaginären Wurzeln muss daher mindestens $e-c+d$ sein.

Zählt man also alle fehlenden Glieder zusammen, rechnet dabei aber in jeder Lücke mit einer ungeraden Anzahl fehlender Glieder bei einem Vorzeichenwechsel eins weniger und bei einer Vorzeichenfolge eins mehr, als Glieder fehlen, so erhält man eine Zahl, die von der Anzahl der imaginären Wurzeln mindestens erreicht werden muss.
\begin{center}\rule{0.28\linewidth}{0.4pt}\end{center}
